\documentclass[12pt]{article}
\usepackage[a4paper,margin=1in]{geometry}
\usepackage[utf8]{inputenc}
\usepackage[T1]{fontenc}
\usepackage{amsmath,amssymb}
\usepackage{textcomp}
\usepackage{parskip}
\begin{document}
\begin{center}
{\Large \textbf{1999 STEP II}}
\end{center}

\subsection*{Page 1}
\begin{flushleft}
STEP II, 1999 2\par
\bigskip\textbf{Section A: Pure Mathematics}\par
\bigskip\textbf{Question 1.} Let x = 10100, y = 10x, z = 10y, and let\par
\[a^1 = x!, a^2 = xy, a^3 = yx, a^4 = zx, a^5 = exyz, a^6 = z 1/y, a^7 = yz/x.\]
\textbf{(i)} Use Stirling's approximation n! \sqrt2\pi nn+ 2 1 e -n, which is valid for large n, to show\par
\approx\par
that log\par
(log\par
\[a^1) 102.\]
\approx\par
\textbf{(ii)} Arrange the seven numbers a1, . . . , a7 in ascending order of magnitude, justifying\par
your result.\par
\bigskip\textbf{Question 2.} Consider the quadratic equation\par
2 2\par
\[nx + 2x\sqrt(pn + q) + rn + s = 0, ( )\]
∗\par
\[where p > 0, p = r and n = 1, 2, 3, . . . .\]
\textbf{(i)} For the case where p = 3, q = 50, r = 2, s = 15, find the set of values of n for which\par
equation ( ) has no real roots.\par
∗\par
\textbf{(ii)} Prove that if p < r and 4q(p r) > s2, then ( ) has no real roots for any value of n.\par
- ∗\par
\textbf{(iii)} If n = 1, p r = 1 and q = s2/8, show that ( ) has real roots if, and only if,\par
- ∗\par
\[s 6 4 2\sqrt\{2\} or s > 4 + 2\sqrt\{2\} .\]
-\par
3 Let\par
\[S (x) =\]
\[ex3\]
dn\par
e\par
\[-x^3\]
.\par
n dxn\par
³ ´\par
\[Show that S^2(x) = 9x4 6x and find S^3(x).\]
-\par
Prove by induction on n that S (x) is a polynomial. By means of your induction argument,\par
n\par
determine the order of this polynomial and the coefficient of the highest power of x.\par
dS\par
\[Show also that if n = 0 for some value a of x, then S\]
n\par
(a)S n+1(a) 6 0.\par
dx\par
\end{flushleft}
\subsection*{Page 2}
\begin{flushleft}
STEP II, 1999 3\par
\bigskip\textbf{Question 4.} By considering the expansions in powers of x of both sides of the identity\par
(1 + x)n(1 + x)n (1 + x) 2n,\par
\equiv\par
show that\par
n 2\par
n 2n\par
\[= ,\]
µs¶ µ n ¶\par
\[Xs=0\]
n n!\par
\[where = .\]
µs¶ s! (n s)!\par
-\par
By considering similar identities, or otherwise, show also that:\par
\textbf{(i)} if n is an even integer, then\par
n 2\par
n n\par
\[( 1)s = ( 1)n/2 ;\]
- µs¶ - µn/2¶\par
\[Xs=0\]
n 2\par
n 2n\par
\textbf{(ii)} 2t = n .\par
µ t ¶ µ n ¶\par
\[Xt=1\]
\bigskip\textbf{Question 5.} Show that if α is a solution of the equation\par
\[5cos x + 12sin x = 7,\]
then either\par
\[35 12\sqrt\{120\}\]
\[cosα = -\]
169\par
or cos α has one other value which you should find.\par
\[Prove carefully that if \pi/2 < α < \pi, then α < 3\pi/4.\]
dy\par
\bigskip\textbf{Question 6.} Find if\par
dx\par
ax + b\par
\[y = . ( )\]
cx + d ∗\par
By using changes of variable of the form ( ), or otherwise, show that\par
∗\par
1 x + 1 1 1 1\par
\[ln dx = ln 3 ln 2 ,\]
\[Z^0 (x + 3)2 µ x + 3 ¶ 6 - 4 - 12\]
and evaluate the integrals\par
\[1 1 x^2 + 3x + 2 1 1 x + 1\]
ln dx and ln dx.\par
\[Z^0 (x + 3)2 µ (x + 3)2 ¶ Z^0 (x + 3)2 µ x + 2 ¶\]
\end{flushleft}
\subsection*{Page 3}
\begin{flushleft}
STEP II, 1999 4\par
\bigskip\textbf{Question 7.} The curve C has equation\par
x\par
\[y = ,\]
\[\sqrt(x^2 2x + a)\]
-\par
\[where the square root is positive. Show that, if a > 1, then C has exactly one stationary\]
point.\par
\[Sketch C when (i) a = 2 and (ii) a = 1.\]
\bigskip\textbf{Question 8.} Prove that\par
n 1 1\par
cos θ cos(n + )θ\par
\[sin kθ = 2 - 2 . ( )\]
2 sin 1 θ ∗\par
\[Xk=0 2\]
\textbf{(i)} Deduce that, when n is large,\par
n\par
\[k\pi 2n\]
sin .\par
\[µ n ¶ \approx \pi\]
\[Xk=0\]
\textbf{(ii)} By differentiating ( ) with respect to θ, or otherwise, show that, when n is large,\par
∗\par
n\par
\[k\pi 1 1\]
2 2\par
k sin + n .\par
\[µ 2n ¶ \approx µ 4 \pi2 ¶\]
\[Xk=0\]
[The approximations, valid for small θ, sin θ θ and cos θ 1 1 θ2 may be assumed.]\par
\approx \approx - 2\par
\end{flushleft}
\subsection*{Page 4}
\begin{flushleft}
STEP II, 1999 5\par
\bigskip\textbf{Section B: Mechanics}\par
\bigskip\textbf{Question 9.} In the Z-universe, a star of mass M suddenly blows up, and the fragments, with various initial\par
speeds, start to move away from the centre of mass G which may be regarded as a fixed point.\par
In the subsequent motion the acceleration of each fragment is directed towards G. Moreover,\par
in accordance with the laws of physics of the Z-universe, there are positive constants k1, k2\par
and R such that when a fragment is at a distance x from G, the magnitude of its acceleration\par
\[is k1x3 if x < R and is k2x -4 if x > R. The initial speed of a fragment is denoted by u.\]
\textbf{(i)} For x < R, write down a differential equation for the speed v, and hence determine\par
\[v in terms of u, k^1 and x for x < R.\]
\textbf{(ii)} Show that if u < a, where 2a2 = k1R4, then the fragment does not reach a distance\par
R from G.\par
\textbf{(iii)} Show that if u > b, where 6b2 = 3k1R4 + 4k2/R3, then from the moment of the\par
explosion the fragment is always moving away from G.\par
\textbf{(iv)} If a < u < b, determine in terms of k2, b and u the maximum distance from G\par
attained by the fragment.\par
\[10 N particles P^1, P^2, P^3, . . ., P\]
N\par
\[with masses m, qm, q2m, . . . , qN-1m, respectively, are at rest\]
at distinct points along a straight line in gravity-free space. The particle P1 is set in motion\par
towards P2 with velocity V and in every subsequent impact the coefficient of restitution is e,\par
\[where 0 < e < 1. Show that after the first impact the velocities of P^1 and P^2 are\]
1 eq 1 + e\par
- V and V,\par
µ 1 + q ¶ µ 1 + q ¶\par
respectively.\par
Show that if q 6 e, then there are exactly N 1 impacts and that if q = e, then the total loss\par
-\par
of kinetic energy after all impacts have occurred is equal to\par
me 1\par
eN-1\par
V\par
.\par
2 -\par
¡ ¢\par
\end{flushleft}
\subsection*{Page 5}
\begin{flushleft}
STEP II, 1999 6\par
\bigskip\textbf{Question 11.} An automated mobile dummy target for gunnery practice is moving anti-clockwise around the\par
circumference of a large circle of radius R in a horizontal plane at a constant angular speed ω.\par
A shell is fired from O, the centre of this circle, with initial speed V and angle of elevation α.\par
Show that if V 2 < gR, then no matter what the value of α, or what vertical plane the shell is\par
fired in, the shell cannot hit the target.\par
Assume now that V 2 > gR and that the shell hits the target, and let β be the angle through\par
which the target rotates between the time at which the shell is fired and the time of impact.\par
Show that β satisfies the equation\par
2 4 2 2 2 2 4\par
\[g β 4ω V β + 4R ω = 0.\]
-\par
Deduce that there are exactly two possible values of β.\par
\[Let β1 and β2 be the possible values of β and let P^1 and P^2 be the corresponding points of\]
impact. By considering the quantities (β2 + β2) and β2β2 , or otherwise, show that the linear\par
1 2 1 2\par
\[distance between P^1 and P^2 is\]
ω\par
\[2R sin \sqrt(V Rg) .\]
g -\par
³ ´\par
\end{flushleft}
\subsection*{Page 6}
\begin{flushleft}
STEP II, 1999 7\par
\bigskip\textbf{Section C: Probability and Statistics}\par
12 It is known that there are three manufacturers A, B, C, who can produce micro chip MB666.\par
The probability that a randomly selected MB666 is produced by A is 2p, and the corresponding\par
probabilities for B and C are p and 1 3p, respectively, where 0 6 p 6 1 . It is also known that\par
- 3\par
\[70\% of MB666 micro chips from A are sound and that the corresponding percentages for B\]
and C are 80\% and 90\%, respectively.\par
\[Find in terms of p, the conditional probability, P(A S), that if a randomly selected MB666\]
|\par
chip is found to be sound then it came from A, and also the conditional probability, P(C S),\par
|\par
that if it is sound then it came from C.\par
\[A quality inspector took a random sample of one MB666 micro chip and found it to be sound.\]
She then traced its place of manufacture to be A, and so estimated p by calculating the value\par
of p that corresponds to the greatest value of P(A S). A second quality inspector also a took\par
|\par
\[random sample of one MB666 chip and found it to be sound. Later he traced its place of\]
manufacture to be C and so estimated p by applying the procedure of his colleague to P(C S).\par
|\par
Determine the values of the two estimates and comment briefly on the results obtained.\par
\bigskip\textbf{Question 13.} A stick is broken at a point, chosen at random, along its length. Find the probability that the\par
ratio, R, of the length of the shorter piece to the length of the longer piece is less than r.\par
Find the probability density function for R, and calculate the mean and variance of R.\par
\bigskip\textbf{Question 14.} You play the following game. You throw a six-sided fair die repeatedly. You may choose\par
to stop after any throw, except that you must stop if you throw a 1. Your score is the\par
number obtained on your last throw. Determine the strategy that you should adopt in order\par
to maximize your expected score, explaining your reasoning carefully.\par
\end{flushleft}
\end{document}
